\documentclass[11pt]{article}
\usepackage[T1]{fontenc}
\usepackage{lmodern}
\usepackage[margin=1.05in]{geometry}
\usepackage{amsmath,amssymb,amsthm,mathtools}
\usepackage{microtype,booktabs,array}
\usepackage{listings}
\usepackage{xcolor}
\usepackage[colorlinks=true,linkcolor=blue!45!black,citecolor=blue!45!black,urlcolor=blue!45!black]{hyperref}
\usepackage[nameinlink,capitalise,noabbrev]{cleveref}
\numberwithin{equation}{section}
\emergencystretch=3em
\newtheorem{theorem}{Theorem}[section]
\newtheorem{proposition}[theorem]{Proposition}
\newtheorem{lemma}[theorem]{Lemma}
\newtheorem{corollary}[theorem]{Corollary}
\newtheorem{conjecture}[theorem]{Conjecture}
\theoremstyle{definition}
\newtheorem{definition}[theorem]{Definition}
\theoremstyle{remark}
\newtheorem{remark}[theorem]{Remark}
\newcommand{\F}{\mathbb F}
\newcommand{\Z}{\mathbb Z}
\DeclareMathOperator{\wt}{wt}
\DeclareMathOperator{\AGL}{AGL}
\lstset{basicstyle=\ttfamily\footnotesize,columns=fullflexible,keepspaces=true,breaklines=true,
  frame=single,framerule=0.3pt,rulecolor=\color{black!30},
  literate={∧}{{$\wedge$}}1 {∀}{{$\forall$}}1 {∃}{{$\exists$}}1 {→}{{$\to$}}2 {≤}{{$\le$}}1
           {⟨}{{$\langle$}}1 {⟩}{{$\rangle$}}1 {↦}{{$\mapsto$}}2}
\title{A balanced 13-variable Boolean function with absolute indicator 112}
\author{Claude (AI model by Anthropic)\\ \small run and submitted by Aleksandr Bersenev}
\date{October 8, 2026\\[3pt]\footnotesize Produced autonomously by Claude (Opus 5.5) agents in Claude Code; run and submitted by Aleksandr Bersenev.\\ \footnotesize Not an official Anthropic publication. Not peer reviewed.}
\hypersetup{pdftitle={A balanced 13-variable Boolean function with absolute indicator 112},pdfauthor={Claude (AI model by Anthropic); submitted by Aleksandr Bersenev}}
\begin{document}
\maketitle
\begin{abstract}
We exhibit an explicit balanced Boolean function $f\colon\F_2^{13}\to\F_2$ of algebraic degree $11$
whose absolute indicator is $\Delta_f=\max_{d\neq 0}\bigl|\sum_{x}(-1)^{f(x)\oplus f(x\oplus d)}\bigr|=112$.
Since $112<2^{\lfloor (13+1)/2\rfloor}=128$, this is a counterexample to the conjecture of Zhang and
Zheng (1995) in the case $n=13$; to our knowledge, for odd $n$ the conjecture had previously been
disproved only for $n\in\{9,11,15,21\}$. The function is invariant under the affine group $\AGL(1,13)$ permuting the
coordinates and was found by a local search in the $2^{74}$-element class of such functions.
Balancedness, the degree bound $\deg f\ge 3$ and $\Delta_f\le 112$ are certified in Lean~4 (core
library only), with all computations checked by the kernel and only the standard axioms.
\end{abstract}
\tableofcontents

\section{Introduction}\label{sec:intro}

For a Boolean function $f\colon\F_2^n\to\F_2$ the \emph{autocorrelation} at $d\in\F_2^n$ is
$r_f(d)=\sum_{x\in\F_2^n}(-1)^{f(x)\oplus f(x\oplus d)}$, and the \emph{absolute indicator} is
$\Delta_f=\max_{d\neq0}|r_f(d)|$. Together with the sum-of-squares indicator, $\Delta_f$ was introduced by
Zhang and Zheng~\cite{ZZ95} as a measure of global avalanche characteristics (GAC); a small $\Delta_f$
means that every nontrivial derivative $x\mapsto f(x)\oplus f(x\oplus d)$ is close to balanced. Bent
functions have $\Delta_f=0$ but are never balanced, and the problem of how small $\Delta_f$ can be for a
balanced function has been studied for three decades. The original statement~\cite[Conjecture~1,
p.~329]{ZZ95} reads: \emph{``Let $f$ be a balanced function on $V_n$, whose algebraic degree is at least
three. Then $\Delta_f\ge2^{\frac12(n+1)}$.''} In the recent literature (we follow Tang and
Maitra~\cite[Conjecture~1]{TM16}) the exponent is written with a floor, which gives a meaningful statement
for even $n$ as well; for odd $n$, in particular for $n=13$, the two forms coincide:

\begin{conjecture}[Zhang--Zheng]\label{conj:ZZ}
The absolute indicator of any balanced Boolean function $f$ on $n$ variables of algebraic degree at
least $3$ satisfies $\Delta_f\ge 2^{\lfloor (n+1)/2\rfloor}$.
\end{conjecture}

\paragraph{Prior counterexamples.} For odd $n$ the conjecture was disproved for $n=15$ by
modifications of Patterson--Wiedemann functions (Maitra--Sarkar~\cite{MS02}, a balanced function with
$\Delta_f=216<256$), for $n=21$ by Patterson--Wiedemann-type functions (Gangopadhyay--Keskar--Maitra~\cite{GKM06}; the
absolute indicator computed there was corrected, and twenty further $21$-variable examples were given, by
Kavut~\cite{Kav16}, see also~\cite{Tang18}), and for $n=9$ ($\Delta_f=24$) and
$n=11$ ($\Delta_f=56$) by Kavut, Maitra and Y\"ucel through a steepest-descent-like search among
rotation-symmetric functions~\cite{KMY06,KMY07}. For even $n$ the conjecture reads $\Delta_f\ge2^{n/2}$
and was disproved for $n=14$ by Burnett, Millan, Dawson and Clark~\cite{BMDC04}, for $n=10$
in~\cite{KMY07}, for $n\equiv 2\pmod 4$, $n\ge 46$ by the construction of Tang and Maitra~\cite{TM16}, and
subsequently for all even $n\ge 12$ by constructions and searches~\cite{KMT19,ZPW21,OKK23}.
The list ``$n=9,11,15,21$'' for odd $n$ is given as the state of the art in the introduction and in the
concluding open problems of~\cite{TM16}, in~\cite[Sec.~1]{KMY07}, and in the survey talk~\cite{Tang18}
(slide ``Some counterexamples on Zhang--Zheng conjecture'', listing $n=9,10,11,14,15,21$), and the same
odd list $n\in\{9,11,15,21\}$ is repeated in the 2019 survey of known results by Canteaut, K\"olsch and
Wiemer~\cite[Sec.~5]{CKW19}; the later papers~\cite{KMT19,ZPW21,OKK23} treat even $n$ only. The case
$n=13$ is the smallest odd case left open. For $n=13$ the conjectured value $128$ is attained, e.g., by the
construction of Zhang and Zheng themselves on $V_{2k+1}$~\cite[Theorem~17]{ZZ95}, while the balanced
$13$-variable functions of record nonlinearity $4036$ of Kavut and Y\"ucel have $\Delta_f=536$~\cite[Appendix~B]{KY07};
we are not aware of any balanced $13$-variable function with $\Delta_f<128$ in the literature.

\paragraph{Result.} We settle the case $n=13$.

\begin{theorem}\label{thm:main}
Let $f$ be the Boolean function on $13$ variables whose truth table is given in Appendix~\ref{app:hex}
(equivalently, the union of the $36$ $\AGL(1,13)$-orbits listed in Appendix~\ref{app:orbits}). Then $f$ is
balanced, $\deg f=11$, and $|r_f(d)|\le 112$ for every $d\neq 0$, with equality for $312$ values of $d$.
In particular $\Delta_f=112<128=2^{\lfloor(13+1)/2\rfloor}$, so Conjecture~\ref{conj:ZZ} fails for $n=13$.
\end{theorem}

The parts ``balanced'', ``$\deg f\ge 3$'' and ``$\Delta_f\le 112$'' are certified in Lean
(\cref{sec:lean}); the exact degree $11$, the equality count and the nonlinearity $3980$ quoted below are
checked by an independent script only. To our knowledge, as of October 8, 2026, no balanced
$13$-variable Boolean function with $\Delta_f<128$ had been reported.

\section{Definitions and elementary facts}\label{sec:defs}

We identify $x=(x_0,\dots,x_{12})\in\F_2^{13}$ with the integer $\sum_i x_i2^i\in\{0,\dots,2^{13}-1\}$;
then $x\oplus d$ is the bitwise exclusive or. A function $f$ is given by the integer
$\mathrm{tt}=\sum_x f(x)2^x<2^{8192}$ (its truth table). It is \emph{balanced} if $\wt(f)=|\{x:f(x)=1\}|=2^{12}$.
The algebraic normal form is $f(x)=\bigoplus_{u}a_u\prod_{i:u_i=1}x_i$ with Möbius coefficients
$a_u=\bigoplus_{x\preceq u}f(x)$, where $x\preceq u$ means $x\wedge u=x$ bitwise, and
$\deg f=\max\{\wt(u):a_u=1\}$. The Walsh transform is $W_f(a)=\sum_x(-1)^{f(x)\oplus a\cdot x}$, the
nonlinearity is $2^{n-1}-\tfrac12\max_a|W_f(a)|$.

\begin{lemma}[Wiener--Khinchin]\label{lem:WK}
$r_f(d)=2^{-n}\sum_{a}W_f(a)^2(-1)^{a\cdot d}$ for all $d$, and
$\sum_{a}\bigl(W_f(a)^2-2^n\bigr)^2=2^n\sum_{d\neq0}r_f(d)^2$.
\end{lemma}
\begin{proof}
$\sum_a W_f(a)^2(-1)^{a\cdot d}=\sum_{x,y}(-1)^{f(x)\oplus f(y)}\sum_a(-1)^{a\cdot(x\oplus y\oplus d)}
=2^n\sum_x(-1)^{f(x)\oplus f(x\oplus d)}$. By Parseval for the pair $(r_f,W_f^2)$ this gives
$\sum_aW_f(a)^4=2^n\sum_d r_f(d)^2$, and with $\sum_aW_f(a)^2=2^{2n}$, $r_f(0)=2^n$ the second identity
follows by expanding the square.
\end{proof}

\begin{lemma}\label{lem:mod8}
If $n\ge3$ and $\wt(f)$ is even (in particular if $f$ is balanced), then $r_f(d)\equiv0\pmod 8$ for all $d$.
\end{lemma}
\begin{proof}
For $d\neq0$ let $S=\{x:f(x)=1\}$. Then $r_f(d)=2^n-2\,\wt(D_df)$ with
$\wt(D_df)=|S\,\triangle\,(S\oplus d)|=2|S|-2|S\cap(S\oplus d)|$. The set $S\cap(S\oplus d)$ is invariant
under the fixed-point-free involution $x\mapsto x\oplus d$, so its size is even; hence
$\wt(D_df)\equiv 2|S|\equiv0\pmod 4$ and $r_f(d)\equiv 2^n\equiv0\pmod 8$.
\end{proof}

Thus for balanced $f$ on $13$ variables the first admissible value below the conjectured bound $128$ is
$120$. We remark also that concatenating two $12$-variable bent functions $g_0\Vert g_1$ always gives
$\Delta\ge128$: here $r(d,1)=2\sum_x(-1)^{g_0(x)\oplus g_1(x\oplus d)}=2W_{\tilde g_0\oplus\tilde g_1}(d)$
(with $\tilde g_i$ the dual bent functions), and $\max_d|W_h(d)|\ge 2^6$ for every $12$-variable $h$.

\section{Construction and search}\label{sec:search}

\paragraph{Symmetry class.} Let $G=\AGL(1,13)=\{i\mapsto ai+b \bmod 13: a\in\Z_{13}^\times,\ b\in\Z_{13}\}$
(order $156$) act on $\F_2^{13}$ by permuting coordinates, $(\pi x)_{\pi(i)}=x_i$. It is generated by the
cyclic shift $i\mapsto i+1$ and the map $i\mapsto 2i$ (as $2$ is a primitive root mod $13$); it contains the
dihedral group, so $G$-invariant functions are in particular rotation- and dihedral-symmetric. The group
has $74$ orbits on $\F_2^{13}$ (sizes $1,13,26,39,52,78,156$), so there are $2^{74}$ invariant functions.
This class appears already in~\cite{KY10}, where Kavut and Y\"ucel used it (as dihedral-symmetric
functions with the extra permutation $x_i\mapsto x_{2i}$) to search for highly nonlinear $13$-variable
functions.

\paragraph{Orbit formulas.} Let $O_1,\dots,O_{74}$ be the orbits with representatives $u_j\in O_j$ and put
$C_{ij}=\sum_{x\in O_i}(-1)^{u_j\cdot x}$. Since permutation matrices are orthogonal,
$W_f(\pi a)=W_f(a)$ and $r_f(\pi d)=r_f(d)$ for $G$-invariant $f$. Writing $s_i=(-1)^{f(O_i)}$ we get
\[
W_f(u_j)=\sum_i C_{ij}s_i,\qquad r_f(u_j)=2^{-13}\sum_i W_f(u_i)^2\,C_{ij}
\]
(the second formula is Lemma~\ref{lem:WK} grouped by orbits). Hence the full Walsh spectrum and autocorrelation
of an invariant function are computed from a $74\times 74$ integer matrix.

\paragraph{Search.} We used the cost function of Kavut and Y\"ucel~\cite{KY03,KMY07},
$\mathrm{cost}(f)=\sum_a(W_f(a)^2-2^{13})^2=2^{13}\sum_{d\ne0}r_f(d)^2$ (Lemma~\ref{lem:WK}), with
moves that complement $f$ on one orbit, steepest descent with a tabu list of length $5$, $3000$ steps
per run and random restarts; after each step $\Delta_f$ and the presence of a zero Walsh coefficient were
recorded. (A zero Walsh coefficient $W_f(a)=0$ means that $f\oplus a\cdot x$ is balanced with the same
$\Delta$; the function below is itself balanced.) Each step costs $O(74^2)$ operations, so thousands of
restarts take seconds. In a run with seed 41 the first function with a Walsh zero and $\Delta\le 120$ appeared after 259
restarts (it is balanced, with $\Delta=120$ and nonlinearity $3984$); the function of \cref{thm:main}
(with $\Delta=112$) was found after $1015$ restarts with the seed reported below, in about $40$ seconds.
Short exploratory runs in the larger classes generated by the shift and $i\mapsto ai$ for $a$ of order
$2,3,4,6$ ($380$, $232$, $198$, $140$ orbits) did not get below $\Delta=136$; a search for
$\Delta\le104$ in $\AGL(1,13)$ found further functions with $\Delta=112$ but none with $\Delta\le 104$.

\paragraph{Reproduction.} With the files in \texttt{data/search/}:
\begin{lstlisting}
python3 t1.py && python3 export.py agl13          # orbits of AGL(1,13), 74x74 matrix
gcc -O3 -march=native -o acls4 acls4.c -lm
./acls4 agl13.bin 13 112 3000 77 100000 5   # class n target steps seed restarts tabu
\end{lstlisting}
It prints a line beginning \texttt{r1015 it2098 ZERO Delta=112 maxW=232 s=}, followed by a $74$-bit string
giving the value of $f$ on the orbits ordered by their least element; this string is stored in
\texttt{data/search/orbit\_string\_agl13.txt} and is equivalent to Appendix~\ref{app:orbits}.

\paragraph{The certificate.} The truth table is the hexadecimal integer in Appendix~\ref{app:hex}. The file
\texttt{data/truth\_table.hex} contains its $2048$ lower-case hex digits followed by a newline; its SHA-256
hash is
\begin{center}\small
\texttt{0b4a77a6d9278877e3f206da6bf4c185bbe29ed76f50726e30843e5a97dba0c8}
\end{center}
(the same digits are the literal \texttt{tt} in \texttt{lean/Main.lean} and Appendix~\ref{app:hex}).
Bit $x$ of the integer is $f(x)$. The script \texttt{data/verify.py} recomputes, by the naive $O(4^{13})$ double sum, the
weight $4096$, all $8191$ values $r_f(d)$, the Walsh spectrum, the algebraic normal form and the
$G$-invariance. Its output: $\deg f=11$ (and $a_{\{0,1,2\}}=1$), $\max_a|W_f(a)|=232$, i.e.\
nonlinearity $3980$ (script-checked only), and the distribution of $|r_f(d)|$, $d\neq0$:

\begin{center}\small
\begin{tabular}{l*{15}{r}}
\toprule
$|r_f(d)|$ & 0 & 8 & 16 & 24 & 32 & 40 & 48 & 56 & 64 & 72 & 80 & 88 & 96 & 104 & 112\\
count & 156 & 416 & 780 & 1392 & 650 & 923 & 416 & 780 & 156 & 585 & 286 & 936 & 247 & 156 & 312\\
\bottomrule
\end{tabular}
\end{center}
All counts except the one for $24$ are multiples of $13$, as they must be for a function invariant under
the cyclic shift: the only shift-fixed nonzero vector is the all-ones vector $\mathbf 1$, and
$r_f(\mathbf 1)=24$. The ANF of $f$ has $4485$ monomials.

\begin{proof}[Proof of \cref{thm:main}]
The statement concerns one explicit function, so the proof is the finite computation; it was carried
out independently by the script above and, for balancedness, $\deg f\ge3$ and $\Delta_f\le112$, by the Lean
kernel as described in \cref{sec:lean}.
\end{proof}

\section{Formal verification}\label{sec:lean}

The file \texttt{lean/Main.lean} ($374$ lines) is checked by Lean~4.34.1~\cite{Lean4} using only the
core library (no Mathlib). There is no \texttt{sorry}, no \texttt{axiom} and no \texttt{native\_decide};
all evaluations use \texttt{decide +kernel}.

\subsection{Statement}
The definitions used by the statement are, verbatim (\texttt{tt} is the $2048$-digit hexadecimal literal of
Appendix~\ref{app:hex}, abbreviated here):
\begin{lstlisting}
def tt : Nat := 0xffeafdcdeba3a0e7...4f1713477dfe4eb75bfea7de3e68
def f (x : Nat) : Bool := tt.testBit x
def sgn (b : Bool) : Int := if b then -1 else 1
def autocorr (d : Nat) : Int := ((List.range (2^13)).map (fun x => sgn (f x ^^ f (x ^^^ d)))).sum
def weight : Nat := ((List.range (2^13)).filter f).length
def anfCoeff (u : Nat) : Bool :=
  ((List.range (2^13)).filter (fun x => x &&& u == x)).foldl (fun b x => b ^^ f x) false
def monoDeg (u : Nat) : Nat := ((List.range 13).filter (fun i => u.testBit i)).length

def MainStatement : Prop :=
  weight = 2^12 ∧
  (∃ u, u < 2^13 ∧ 3 ≤ monoDeg u ∧ anfCoeff u = true) ∧
  (∀ d, 0 < d → d < 2^13 → (autocorr d).natAbs ≤ 112)
\end{lstlisting}
The certified theorems are
\begin{lstlisting}
theorem main :
    weight = 2^12 ∧
    (∃ u, u < 2^13 ∧ 3 ≤ monoDeg u ∧ anfCoeff u = true) ∧
    (∀ d, 0 < d → d < 2^13 → (autocorr d).natAbs ≤ 112)
theorem main_statement : MainStatement := main
theorem zhang_zheng_fails_n13 :
    ∃ g : Nat → Bool,
      ((List.range (2^13)).filter g).length = 2^12 ∧
      (∃ u, u < 2^13 ∧ 3 ≤ monoDeg u ∧
        ((List.range (2^13)).filter (fun x => x &&& u == x)).foldl (fun b x => b ^^ g x) false = true) ∧
      (∀ d, 0 < d → d < 2^13 →
        (((List.range (2^13)).map (fun x => sgn (g x ^^ g (x ^^^ d)))).sum).natAbs < 2^((13+1)/2))
\end{lstlisting}
\texttt{MainStatement} is placed in the file immediately after the definitions; the identity
\texttt{main\_statement} guarantees that it is literally the statement proved at the end.

\subsection{Statement audit}
\begin{itemize}
\item \emph{Encoding.} Inputs are naturals $x<2^{13}$ whose bit $i$ is $x_i$; \texttt{x \textasciicircum{}\textasciicircum{}\textasciicircum{} d} is
bitwise xor, i.e.\ translation by $d$ in $\F_2^{13}$; \texttt{List.range (2\textasciicircum{}13)} enumerates $\F_2^{13}$
exactly once. Only bits $x<2^{13}$ of \texttt{tt} are ever read (and $\mathrm{tt}<2^{8192}$ anyway).
\item \emph{Autocorrelation.} \texttt{autocorr d} is literally $\sum_{x}(-1)^{f(x)\oplus f(x\oplus d)}$
(\texttt{sgn true = -1}); the quantifier ranges over all $d$ with $0<d<2^{13}$, i.e.\ all nonzero shifts,
and \texttt{natAbs} is the absolute value. In \texttt{zhang\_zheng\_fails\_n13} the inequality is strict and
\texttt{(13+1)/2} is natural-number division, equal to $7=\lfloor 14/2\rfloor$.
\item \emph{Balanced.} \texttt{weight = 2\textasciicircum{}12} says exactly half of the $8192$ inputs are mapped to $1$.
\item \emph{Degree.} \texttt{anfCoeff u} is the Möbius coefficient $a_u=\bigoplus_{x\preceq u}f(x)$ (with
$x\preceq u\iff x\wedge u=x$), the coefficient of $\prod_{i\in u}x_i$ in the ANF, and \texttt{monoDeg u}
is the number of variables of that monomial; so the second conjunct is exactly $\deg f\ge3$
(the witness is $u=7$, the monomial $x_0x_1x_2$).
\item The auxiliary functions below never occur in the statements.
\end{itemize}

\subsection{From the naive definitions to a kernel-feasible computation}
A literal evaluation of \texttt{autocorr} for all $d$ needs $2^{26}$ list steps, far beyond what the kernel
can do (a test with $2^{18}$ steps already exhausted 8\,GB). The proof therefore evaluates a fast version
built from \texttt{Nat} bitwise operations, which the Lean kernel executes with GMP, and proves it equal to
the naive definitions by general lemmas:
\begin{enumerate}
\item \texttt{sum\_sgn}: for any list $L$ and predicate $p$, \texttt{(L.map (fun x => sgn (p x))).sum =
L.length - 2 * (L.filter p).length}. Hence \texttt{autocorr d} $=8192-2N_d$ with
$N_d=|\{x:f(x)\ne f(x\oplus d)\}|$.
\item \emph{Translation by block swaps.} With masks
\texttt{M i} (ones exactly at positions $x<2^{13}$ whose bit $i$ is $0$) and
\texttt{swapBlk i U := ((U >{}>{}> 2\textasciicircum{}i) \&\&\& M i) ||| ((U \&\&\& M i) <{}<{}< 2\textasciicircum{}i)}, lemma
\texttt{swapBlk\_testBit} proves \texttt{(swapBlk i U).testBit x = U.testBit (x \textasciicircum{}\textasciicircum{}\textasciicircum{} 2\textasciicircum{}i)} for all
$U$, $i<13$, $x<2^{13}$. Its proof is a case analysis on bit $i$ of $x$ using the core lemmas on
\texttt{testBit} of \texttt{\&\&\&}, \texttt{|||}, \texttt{>{}>{}>}, \texttt{<{}<{}<}, together with $13\cdot 8192$
elementary facts about the masks and about $x\oplus 2^i$ (\texttt{factOK}), checked by \texttt{decide +kernel}
in $13$ separate theorems. Composing the swaps for the bits of $d$ gives \texttt{perm d T} with
\texttt{perm\_testBit : (perm d T).testBit x = T.testBit (x \textasciicircum{}\textasciicircum{}\textasciicircum{} d)} for $d,x<2^{13}$ (induction over
the bits, plus a kernel check that the accumulated xor mask equals $d$ for all $d<2^{13}$).
\item \emph{Population count modulo $2^{14}-1$.} Let $D_q=\sum_{j<q}2^{14j}$ and
$Y(V,q)=\sum_{t=0}^{13}\bigl((V\gg t)\wedge D_q\bigr)$. Lemma \texttt{and\_D\_succ}
($A\wedge(2^{14}E+1)=2^{14}((A\gg 14)\wedge E)+(A\bmod 2)$) gives \texttt{Y\_succ}:
$Y(V,q+1)=2^{14}Y(V\gg14,q)+\#\{t<14: V_t=1\}$, and since $2^{14}\equiv1\pmod{2^{14}-1}$, induction gives
\texttt{Y\_mod}: $Y(V,q)\equiv \#\{x<14q:V_x=1\}\pmod{16383}$. With $q=586$ ($14q=8204$) and the low
$8192$ bits masked, \texttt{popc\_eq : popc V = cntBits V 8192} where
\texttt{popc V := Y (V \&\&\& (2\textasciicircum{}8192-1)) 586 \% 16383} and \texttt{cntBits V m} is the naive count of set
bits below $m$.
\item \texttt{autocorr\_eq}: for $d<2^{13}$, \texttt{autocorr d = 8192 - 2 * popc (tt \textasciicircum{}\textasciicircum{}\textasciicircum{} perm d tt)},
from 1--3.
\item \texttt{chk0},\dots,\texttt{chk7}: the kernel verifies $4040\le$ \texttt{popc (tt \textasciicircum{}\textasciicircum{}\textasciicircum{} perm d tt)}
$\le4152$ for all $0<d<2^{13}$, in eight chunks of $1024$ shifts; each shift costs roughly a hundred
operations on $8192$-bit integers. Together with 4 this is $|r_f(d)|\le 112$.
\item \texttt{weight\_eq}: \texttt{weight = popc tt = 4096}; and \texttt{anfCoeff 7 = true},
\texttt{monoDeg 7 = 3}, both by \texttt{decide +kernel}.
\end{enumerate}

\subsection{Axioms and resources}
The output of \texttt{lean Main.lean}, run in \texttt{lean/} with Lean~4.34.1 through a small timing
wrapper (\texttt{leanchk}, which appends the last line below; the wrapper itself is not distributed) and
recorded in \texttt{lean/lean\_output.txt}, is
\begin{lstlisting}
'A12.main' depends on axioms: [propext, Classical.choice, Quot.sound]
'A12.main_statement' depends on axioms: [propext, Classical.choice, Quot.sound]
'A12.zhang_zheng_fails_n13' depends on axioms: [propext, Classical.choice, Quot.sound]
leanchk: wall=124.31s maxrss=2369432kB exit=0
\end{lstlisting}
i.e.\ only the three standard axioms; in particular no \texttt{Lean.ofReduceBool}. The option
\texttt{Elab.async false} is set so that the kernel checks run sequentially within the memory limit.
An independent re-run of \texttt{lean Main.lean} with a freshly installed Lean~4.34.1 core toolchain
(review of October 8, 2026) reproduced the three axiom lines exactly, with exit code $0$, wall time
$124.1$\,s and maximum resident set size $2.36$\,GB; the same review also reproduced the two search runs
of \cref{sec:search} (seeds $77$ and $41$) bit for bit.

\section{Limitations and open questions}\label{sec:limits}
\begin{itemize}
\item Novelty rests on a literature search (IACR ePrint, arXiv including 2025--2026 papers on odd-$n$
Boolean functions, journal pages and survey slides); every source we could check that lists the known
counterexamples~\cite{TM16,Tang18,CKW19} gives $n\in\{9,11,15,21\}$ for odd $n$, and to our knowledge, as
of October 8, 2026, the case $n=13$ had not been settled. The conjecture is quoted from the original
paper~\cite{ZZ95} in \cref{sec:intro}; the floor in the exponent of \cref{conj:ZZ} follows~\cite{TM16}.
\item Only $n=13$ is settled. For odd $n$ the cases $n=17,19$ appear to remain open; our method extends
in principle (the class $\AGL(1,p)$ for prime $p$) but kernel verification of $n=17$ would be about
$2^8$ times more expensive.
\item The minimum of $\Delta_f$ over balanced $13$-variable functions is unknown; we found $112$ but not
$104$ in $\AGL(1,13)$. The nonlinearity $3980$ of our function is below the bent concatenation bound
$4032$; whether $\Delta_f<128$ is compatible with nonlinearity $>4032$ for $n=13$ is open.
\item Degree $11$, nonlinearity $3980$, the $|r_f|$ distribution and the $\AGL(1,13)$-invariance are
script-checked only.
\end{itemize}

\begin{remark}[Negative results on balanced nonlinearity]
Before turning to the absolute indicator we looked for balanced functions beating the best known
balanced nonlinearities $992$ ($n=11$) and $4036$ ($n=13$)~\cite{KY07,Mai07,CDJMP22}, starting from
$h(y)\oplus f(x)$ with $f$ one of the published $9$-variable functions of nonlinearity $242$~\cite{KY10}.
An exhaustive search found that toggling $4$ bits of any of $69$ such $11$-variable bases never gives a
balanced function of nonlinearity $\ge994$, and tabu search over $8$-bit toggles of the $13$-variable
bases (the route to $4036$ in~\cite{KY07}) never approached $4038$. These computations are not certified
and are recorded only as negative evidence.
\end{remark}

\appendix
\section{Truth table}\label{app:hex}
$f(x)$ is bit $x$ (weight $2^x$) of the following $8192$-bit integer, written in hexadecimal with the most
significant digit first, $64$ digits per line (the same digits as \texttt{data/truth\_table.hex} and the
literal \texttt{tt} in \texttt{lean/Main.lean}):
{\small
\begin{verbatim}
ffeafdcdeba3a0e7e8cf885b8914b96efd80b4afd1d436dac5c60734de9228e9
eff6d450ce35c8bee247f7205e69a3d8a137e438517b5b64f3b8861c1cd4a882
e8aafe78e3253354a1ad0f67f4d1dafdb81d206beb7a0d0036e87987995ff3c1
8d171b6ef9255ac436122fce638b6861bb4ade94907d47a516e1f3758cc5815c
f9d18dc9faed2f84b91f1c671b1b6675dc1398f244bf3c2aeb75f743b298feb2
8bd112a30c543d9bf99e6ed844e240510b2cfcc47ed3802b828662affa0be013
c1f2463a538f78ecaed7586236c9a1705e6857480ceee0e8394a859b2c817d47
de8e65d8b2a8c220c3102af3642b98334669f842fa4e7b37c4a0e176840767b5
ef83b213c5b6f5c2fed9bce21cab85708e9317eb12a4693a539b478a2d3d6b73
a6b5170fc2d4be0c20749eea1fe01988ac9b7a62ee7a244fde0882c0eab9cb5d
c0dee64646188c5a55b536604ae7c28bbfc6d6b87ce9e2847570b94970112712
15cf1da4bfe4a0206bede31e8511498f945d9039381c88bfbbd804dfa905035b
a016af1c75295edc330fc4ae6ac0bcb1c8edf66f67d12c095a69a09298532b11
36e969c5666a648545e0b8eda814fdd15b83249d8027928e1da1c0077bb2756f
e3ac85b83c37e695da0c9880a14d4815b15f4301088def0b6821498bc2d45e5f
607c79d7fe857119bb8935b96f8e5e6fe071d914a8132f3d8131116f3d3fdb76
f8fb910b8b0d165af063de6cba36e409beeca7868bf0e91c17a4dddad4266e00
80ec935e133bedcf47099c3068d35fd9375ad3ca643a918958e64fb369df2b1e
dd39cb32576b55afe54ca7609ff944e41d412e75d6aca88d12eaed010283c185
9ca4d3ca2fc8691ce8f82a88583521fbe2bc4080940da145a9c8cad7a09e23b7
b145b6f8ed68613c25794294d5b0379d27738e620e68785025c8ac2be008859f
8eaaa42cf7688b817ae4e982fc59c5347b727a009b9635936f4546075c2f131d
0636f1ee53a288658bebbd21c9044d15699ba8a3b84b06fdd12703427183d4fb
923467a7c2401b875f9103b181918eef8edfe3d55175b2eb8dd60537155b77de
cd51136cc9ea46a16a321c9672fce2a00a5e55bbb1648dfd7cd8a451cab59f16
f594f9f3ba387daa792fe20348b1019362d838968c00c71d8680631b188f0617
5b6cb8963d86f4773d2828cc382180372022f811ca90edb6dcc05675bba3e357
26dad00b197192b79501492fd318d1ed46b38853f515152b2f9fdf493e377dfe
bd5a98a18466ce800ef55a3fec2d8776f6dd44f0d38185158c4620b720954327
da4672ef700a410700c584b6e8ef519f2cc15d477193d4dba50df77137f967fe
7c412aa02fc7a33eefacd1372a070793cacbc0930b268bd729fbd5ed33fd7dbe
e9003b16b2d703358985121f1dfb5ff690564f1713477dfe4eb75bfea7de3e68
\end{verbatim}

}

\section{Orbit description}\label{app:orbits}
Equivalently, $f(x)=1$ if and only if the $\AGL(1,13)$-orbit of $x$ has least element (as an integer) in
\begin{center}\small
\begin{tabular}{p{0.9\textwidth}}
3, 11, 19, 23, 31, 39, 43, 55, 107, 127, 151, 183, 187, 191, 215, 223, 255, 303, 311, 367, 375, 383,
431, 443, 447, 471, 763, 767, 831, 895, 951, 1023, 1791, 2047, 4095, 8191.
\end{tabular}
\end{center}
These $36$ orbits have sizes (in the same order) $78$, $156$, $52$, $156$, $78$, $156$, $156$, $156$, $39$, $78$,
$156$, $156$, $156$, $156$, $156$, $156$, $78$, $156$, $52$, $156$, $156$, $156$, $156$, $26$, $156$, $78$, $156$, $156$,
$156$, $156$, $39$, $78$, $52$, $78$, $13$, $1$, summing to $4096$.

{\small
\bibliographystyle{plain}
\bibliography{references}
}
\end{document}
